\section{Konu 5: Kısmi Türevler, Teğet Düzlemler ve Lineer Yaklaşımlar}

Bu bölümde 22.10 (defter notu), 05.11 ve 08.11 derslerinde tahtada çözülen tüm sorular derste işlenme sırasıyla ve adım
adım çözülmüştür. 08.11 dersinde 22.10 konuları (kısmi türev, Schwarz teoremi, lineerleştirme) ve 05.11 problemleri
tahtada yeniden yazıldığı için aynı soru birden fazla fotoğrafta görünüyorsa tüm fotoğraflar birlikte verilmiştir.
Bu konuya ait tahtada sorusu yazılıp çözülmeden bırakılmış bir çalışma sorusu yoktur; ``Problemler'' başlığıyla verilen
soruların hepsi derste çözülmüştür. Sonda kaynak kitaplardaki tipik alıştırmalara benzer ek alıştırmalar vardır.

%-------------------------------------------------------------------------------------------
\subsection{Kısmi türevler ve yüksek mertebeden kısmi türevler (22.10 dersi)}

\begin{soru}[Soru 1 \foto{IMG-20261022-WA0055}, \foto{IMG-20261108-WA0090}]
$f(x,y)=x^2y^3+3x$ fonksiyonunun $f_x$ ve $f_y$ kısmi türevlerini bulunuz.
\end{soru}
\begin{cevap}[Çözüm (adım adım)]
\textbf{Kural:} Bir değişkene göre kısmi türev alırken diğer değişken sabit (bir sayı gibi) kabul edilir.

\textbf{$f_x$ ($y$ sabit):}
\begin{itemize}
  \item $x^2y^3$: $y^3$ sabit katsayıdır, $\dfrac{d}{dx}(x^2)=2x$ $\Rightarrow$ $2xy^3$.
  \item $3x$: türevi $3$.
\end{itemize}
\[ f_x=2xy^3+3 . \]
\textbf{$f_y$ ($x$ sabit):}
\begin{itemize}
  \item $x^2y^3$: $x^2$ sabit katsayıdır, $\dfrac{d}{dy}(y^3)=3y^2$ $\Rightarrow$ $3x^2y^2$.
  \item $3x$: $y$ içermediği için $y$'ye göre sabittir; türevi $0$.
\end{itemize}
\[ f_y=3x^2y^2 . \]
\textbf{Ek kontrol (Schwarz teoremi):} $f_{xy}=(2xy^3+3)_y=6xy^2$, $f_{yx}=(3x^2y^2)_x=6xy^2$; karma türevler eşittir
(fonksiyon simetrik olmadığı hâlde).
\end{cevap}

\begin{soru}[Soru 2 \foto{IMG-20261022-WA0055}, \foto{IMG-20261022-WA0059}, \foto{IMG-20261108-WA0091}, \foto{IMG-20261108-WA0092}]
$f(x,y)=\dfrac{xy}{x^2+y^2}$ fonksiyonunun 2. mertebeden kısmi türevlerini hesaplayınız.
\end{soru}
\begin{cevap}[Çözüm (adım adım)]
Fonksiyon $(x,y)\ne(0,0)$ için tanımlıdır; tüm hesaplar bu noktalarda geçerlidir. Bölüm kuralı:
$\left(\dfrac uv\right)'=\dfrac{u'v-uv'}{v^2}$.

\textbf{Adım 1 -- $f_x$.} $y$ sabit: $u=xy\Rightarrow u_x=y$; $v=x^2+y^2\Rightarrow v_x=2x$.
\[
f_x=\frac{y(x^2+y^2)-2x\cdot xy}{(x^2+y^2)^2}=\frac{yx^2+y^3-2x^2y}{(x^2+y^2)^2}=\frac{y^3-x^2y}{(x^2+y^2)^2}.
\]
\textbf{Adım 2 -- $f_{xx}$.} $f_x$'i $x$'e göre türevle: $u=y^3-x^2y\Rightarrow u_x=-2xy$;
$v=(x^2+y^2)^2\Rightarrow v_x=2(x^2+y^2)\cdot2x$.
\[
f_{xx}=\frac{-2xy(x^2+y^2)^2-2(x^2+y^2)\cdot2x\,(y^3-x^2y)}{(x^2+y^2)^4}.
\]
Pay ve payda $(x^2+y^2)$ ile sadeleşir:
\[
f_{xx}=\frac{-2xy(x^2+y^2)-4x(y^3-x^2y)}{(x^2+y^2)^3}=\frac{-2x^3y-2xy^3-4xy^3+4x^3y}{(x^2+y^2)^3}
=\frac{2x^3y-6xy^3}{(x^2+y^2)^3}.
\]
\textbf{Adım 3 -- $f_{xy}$.} $f_x$'i $y$'ye göre türevle: $u=y^3-x^2y\Rightarrow u_y=3y^2-x^2$;
$v_y=2(x^2+y^2)\cdot2y$.
\begin{align*}
f_{xy}&=\frac{(3y^2-x^2)(x^2+y^2)^2-2(x^2+y^2)\cdot2y\,(y^3-x^2y)}{(x^2+y^2)^4}
=\frac{(3y^2-x^2)(x^2+y^2)-4y(y^3-x^2y)}{(x^2+y^2)^3}\\
&=\frac{3x^2y^2+3y^4-x^4-x^2y^2-4y^4+4x^2y^2}{(x^2+y^2)^3}
=\boxed{\frac{-x^4-y^4+6x^2y^2}{(x^2+y^2)^3}} .
\end{align*}
\textbf{Adım 4 -- $f_y$.} $x$ sabit: $u_y=x$, $v_y=2y$.
\[
f_y=\frac{x(x^2+y^2)-2y\cdot xy}{(x^2+y^2)^2}=\frac{x^3-xy^2}{(x^2+y^2)^2}.
\]
\textbf{Adım 5 -- $f_{yx}$.} $f_y$'yi $x$'e göre türevle: $u=x^3-xy^2\Rightarrow u_x=3x^2-y^2$; $v_x=2(x^2+y^2)\cdot2x$.
\begin{align*}
f_{yx}&=\frac{(3x^2-y^2)(x^2+y^2)^2-2(x^2+y^2)\cdot2x\,(x^3-xy^2)}{(x^2+y^2)^4}
=\frac{(3x^2-y^2)(x^2+y^2)-4x(x^3-xy^2)}{(x^2+y^2)^3}\\
&=\frac{3x^4+3x^2y^2-x^2y^2-y^4-4x^4+4x^2y^2}{(x^2+y^2)^3}
=\boxed{\frac{-x^4-y^4+6x^2y^2}{(x^2+y^2)^3}} .
\end{align*}
\textbf{Adım 6 -- $f_{yy}$.} $f_y$'yi $y$'ye göre türevle: $u_y=-2xy$, $v_y=2(x^2+y^2)\cdot2y$.
\[
f_{yy}=\frac{-2xy(x^2+y^2)-4y(x^3-xy^2)}{(x^2+y^2)^3}=\frac{-2x^3y-2xy^3-4x^3y+4xy^3}{(x^2+y^2)^3}
=\frac{2xy^3-6x^3y}{(x^2+y^2)^3}.
\]
\textbf{Sonuç:} $f_{xy}(x,y)=f_{yx}(x,y)$. Bu eşitlik, karma türevler $(0,0)$ dışında sürekli olduğu için
\textbf{Schwarz (Clairaut) teoreminin} beklediği sonuçtur: \emph{$f_x$, $f_y$ türevlenebilir ve karma türevler sürekli
ise $f_{xy}=f_{yx}$}.
\end{cevap}

\begin{dikkat}[Tahtadaki notlarla ilgili düzeltmeler]
\begin{itemize}
  \item $f_{xy}$'nin ara satırında ``$-4y^2$'' yazılmıştı (\foto{IMG-20261022-WA0055}, \foto{IMG-20261108-WA0091});
        doğrusu $-4y\cdot y^3=-4y^4$'tür. Defter notundaki $f_{xx}$ ara satırının ilk terimi de ``$2x^3y$'' olarak
        görünüyor; doğrusu $-2x^3y$ (\foto{IMG-20261108-WA0091}'de doğru). Kutulanan sonuçlar doğrudur.
  \item Defter notunda (\foto{IMG-20261022-WA0059}) ``$f_{xy}=f_{yx}$, eğer $f(x,y)=f(y,x)$ ise geçerlidir'' yazılmıştı.
        Bu yanlıştır: karma türevlerin eşitliği simetriye değil, karma türevlerin \textbf{sürekliliğine} bağlıdır
        (Soru 1'deki simetrik olmayan fonksiyonda da $f_{xy}=f_{yx}$).
\end{itemize}
\end{dikkat}

%-------------------------------------------------------------------------------------------
\subsection{Lineerleştirme ve hata sınırı (22.10 dersi)}

\begin{soru}[Soru 3 \foto{IMG-20261022-WA0060}, \foto{IMG-20261108-WA0093}]
$f(x,y)=x^2+3xy+2y^2+4$ fonksiyonunun $(2,1)$ noktasındaki lineerleştirmesini bulunuz.
$B=\{(x,y):|x-2|\le0{,}01,\ |y-1|\le0{,}01\}$ dikdörtgensel bölgesi üzerindeki hatanın bir üst sınırını hesaplayınız.
\end{soru}
\begin{cevap}[Çözüm (adım adım)]
\textbf{Kullanılan formüller:}
$L(x,y)=f(x_0,y_0)+f_x(x_0,y_0)(x-x_0)+f_y(x_0,y_0)(y-y_0)$ ve
$|E(x,y)|\le\dfrac12M\big(|x-x_0|+|y-y_0|\big)^2$, burada $M$, $B$ üzerinde $|f_{xx}|,|f_{xy}|,|f_{yy}|$'nin üst sınırı.

\textbf{Adım 1 -- Fonksiyon değeri:} $f(2,1)=2^2+3\cdot2\cdot1+2\cdot1^2+4=4+6+2+4=16$.

\textbf{Adım 2 -- Birinci kısmi türevler:}
$f_x=2x+3y\Rightarrow f_x(2,1)=4+3=7$;\quad $f_y=3x+4y\Rightarrow f_y(2,1)=6+4=10$.

\textbf{Adım 3 -- Lineerleştirme:}
\[ L(x,y)=16+7(x-2)+10(y-1)=16+7x-14+10y-10=7x+10y-8 . \]
\textbf{Adım 4 -- İkinci türevler ve $M$:} $f_{xx}=2$, $f_{xy}=3$, $f_{yy}=4$. Sabit oldukları için $B$'nin her yerinde
aynıdır; $M=\max\{2,3,4\}=4$.

\textbf{Adım 5 -- Hata sınırı:} $B$'de $|x-2|\le0{,}01$ ve $|y-1|\le0{,}01$:
\[
|E(x,y)|\le\frac12\cdot4\,\big(|x-2|+|y-1|\big)^2\le2\,(0{,}01+0{,}01)^2=2\cdot0{,}0004=0{,}0008 .
\]
\textbf{Yorum:} $B$ içinde $f$ yerine $7x+10y-8$ kullanılırsa hata en fazla $0{,}0008$ olur. (Bu $f$ ikinci derece polinom
olduğundan hatanın tam ifadesi $E=(x-2)^2+3(x-2)(y-1)+2(y-1)^2$'dir; $B$'deki en büyük değeri $0{,}0006$'dır, yani bulunan
sınır geçerlidir.)
\end{cevap}

%-------------------------------------------------------------------------------------------
\subsection{Tam diferansiyel ve diferansiyelle yaklaşık hesap -- Problemler (05.11 dersi)}

\begin{soru}[Soru 4 \foto{IMG-20261105-WA0034}, \foto{IMG-20261108-WA0082} (Problem 1a)]
$f(x,y)=x^3+5xy^2+3xy+4x-3y+5$ fonksiyonunun birinci ve ikinci mertebeden tam diferansiyellerini bulunuz.
\end{soru}
\begin{cevap}[Çözüm (adım adım)]
\textbf{Formüller:} $df=f_x\,dx+f_y\,dy$,\quad $d^2f=f_{xx}\,dx^2+2f_{xy}\,dx\,dy+f_{yy}\,dy^2$.

\textbf{Birinci türevler:} $f_x=3x^2+5y^2+3y+4$ ($-3y$ ve $5$'in $x$'e göre türevi $0$),\quad
$f_y=10xy+3x-3$ ($x^3$, $4x$, $5$'in $y$'ye göre türevi $0$).
\[ df=(3x^2+5y^2+3y+4)\,dx+(10xy+3x-3)\,dy . \]
\textbf{İkinci türevler:} $f_{xx}=6x$,\quad $f_{xy}=(f_x)_y=10y+3$ \ (kontrol: $f_{yx}=(10xy+3x-3)_x=10y+3$ $\checkmark$),\quad
$f_{yy}=10x$.
\[ d^2f=6x\,dx^2+2(10y+3)\,dx\,dy+10x\,dy^2 . \]
\end{cevap}

\begin{soru}[Soru 5 \foto{IMG-20261105-WA0034}, \foto{IMG-20261105-WA0035}, \foto{IMG-20261108-WA0082}, \foto{IMG-20261108-WA0083} (Problem 1b)]
$w=f(x,y,z)=3x^2yz^2-3xy+2yz^2+xyz$ fonksiyonunun birinci ve ikinci mertebeden tam diferansiyellerini bulunuz.
\end{soru}
\begin{cevap}[Çözüm (adım adım)]
\textbf{Formüller (üç değişken):} $dw=w_x\,dx+w_y\,dy+w_z\,dz$,
\[ d^2w=w_{xx}dx^2+w_{yy}dy^2+w_{zz}dz^2+2w_{xy}dx\,dy+2w_{xz}dx\,dz+2w_{yz}dy\,dz . \]
\textbf{Birinci türevler} (her seferinde diğer iki değişken sabit):
\begin{align*}
w_x&=6xyz^2-3y+yz, & w_y&=3x^2z^2-3x+2z^2+xz, & w_z&=6x^2yz+4yz+xy .
\end{align*}
\[ dw=(6xyz^2-3y+yz)\,dx+(3x^2z^2-3x+2z^2+xz)\,dy+(6x^2yz+4yz+xy)\,dz . \]
\textbf{İkinci türevler:}
\begin{align*}
w_{xx}&=(w_x)_x=6yz^2, & w_{yy}&=(w_y)_y=0, & w_{zz}&=(w_z)_z=6x^2y+4y,\\
w_{xy}&=(w_x)_y=6xz^2-3+z, & w_{xz}&=(w_x)_z=12xyz+y, & w_{yz}&=(w_y)_z=6x^2z+4z+x .
\end{align*}
(Kontrol: $(w_y)_x=6xz^2-3+z=w_{xy}$, $(w_z)_x=12xyz+y=w_{xz}$, $(w_z)_y=6x^2z+4z+x=w_{yz}$ $\checkmark$.)

$w_{yy}=0$ olduğundan $dy^2$ terimi yoktur:
\begin{align*}
d^2w={}&6yz^2\,dx^2+(6x^2y+4y)\,dz^2+2(6xz^2+z-3)\,dx\,dy\\
&+2(12xyz+y)\,dx\,dz+2(6x^2z+x+4z)\,dy\,dz .
\end{align*}
\end{cevap}

\begin{soru}[Soru 6 \foto{IMG-20261105-WA0035}, \foto{IMG-20261108-WA0083}, \foto{IMG-20261108-WA0084} (Problem 2)]
Bir dik dairesel silindirin yarıçapı 10 cm'den 9,9 cm'ye azalıyor, yüksekliği 30 cm'den 30,3 cm'ye çıkıyor. Hacimdeki
değişimi diferansiyel ile yaklaşık olarak hesaplayınız.
\end{soru}
\begin{cevap}[Çözüm (adım adım)]
\textbf{Adım 1 -- Fonksiyon:} $V=\pi r^2h$.

\textbf{Adım 2 -- Tam diferansiyel:} $V_r=2\pi rh$, $V_h=\pi r^2$ $\Rightarrow$ $dV=2\pi rh\,dr+\pi r^2\,dh$.

\textbf{Adım 3 -- Değerler:} $r=10$, $h=30$; $dr=9{,}9-10=-0{,}1$ (azalma, negatif), $dh=30{,}3-30=0{,}3$.

\textbf{Adım 4 -- Hesap:}
\[
dV=2\pi\cdot10\cdot30\cdot(-0{,}1)+\pi\cdot100\cdot0{,}3=-60\pi+30\pi=-30\pi\ \text{cm}^3\approx-94{,}25\ \text{cm}^3 .
\]
\textbf{Adım 5 -- Yeni hacim:} $V=\pi\cdot100\cdot30=3000\pi$ $\Rightarrow$ yeni hacim $\cong3000\pi-30\pi=2970\pi$ cm$^3$.

Hacim yaklaşık $30\pi$ cm$^3$ azalır. (Kontrol: gerçek yeni hacim $\pi\cdot9{,}9^2\cdot30{,}3\approx2969{,}70\pi$ cm$^3$,
gerçek değişim $\approx-30{,}30\pi$ cm$^3$.)
\end{cevap}

\begin{soru}[Soru 7 \foto{IMG-20261105-WA0036}, \foto{IMG-20261108-WA0085} (Problem 3)]
Yarıçapı $r=1$, yüksekliği $h=5$ olan bir silindirin yarıçapı $0{,}03$ arttırılıp, yüksekliği $0{,}1$ azaltılırsa
hacimdeki yaklaşık değişim ve yüzde değişim nedir?
\end{soru}
\begin{cevap}[Çözüm (adım adım)]
\textbf{Adım 1:} $dV=2\pi rh\,dr+\pi r^2\,dh$; $r=1$, $h=5$, $dr=0{,}03$, $dh=-0{,}1$.

\textbf{Adım 2:} $dV=2\pi\cdot1\cdot5\cdot0{,}03+\pi\cdot1\cdot(-0{,}1)=0{,}3\pi-0{,}1\pi=0{,}2\pi\approx0{,}63$.

\textbf{Adım 3 -- Bağıl ve yüzde değişim:} Başlangıç hacmi $V=\pi\cdot1^2\cdot5=5\pi$.
\[ \frac{dV}{V}=\frac{0{,}2\pi}{5\pi}=0{,}04\quad\Rightarrow\quad\text{hacim yaklaşık \%4 artar.} \]
\textbf{Alternatif (logaritmik diferansiyel):} $\ln V=\ln\pi+2\ln r+\ln h$ $\Rightarrow$
$\dfrac{dV}{V}=2\dfrac{dr}{r}+\dfrac{dh}{h}=2(0{,}03)+\dfrac{-0{,}1}{5}=0{,}06-0{,}02=0{,}04$.

(Kontrol: gerçek değişim $\pi(1{,}03^2\cdot4{,}9-5)\approx0{,}198\pi$, yani yaklaşık \%3,97.)
\end{cevap}
Not: \foto{IMG-20261108-WA0085}'in sol kenarında bu sorunun yalnızca kesik bir parçası görünür; soru metni ve çözüm
\foto{IMG-20261105-WA0036}'dan alınmıştır.

%-------------------------------------------------------------------------------------------
\subsection{Lineerleştirme ve hata sınırı -- Problemler (05.11 dersi)}

\begin{soru}[Soru 8 \foto{IMG-20261105-WA0037}, \foto{IMG-20261108-WA0085} (Problem 4)]
$f(x,y)=e^{2y-x}$ fonksiyonunun $P(0,0)$ ve $Q(1,2)$ noktalarındaki lineerleştirmelerini bulunuz.
\end{soru}
\begin{cevap}[Çözüm (adım adım)]
\textbf{Kısmi türevler} (zincir kuralı; iç fonksiyon $u=2y-x$, $u_x=-1$, $u_y=2$):
$f_x=-e^{2y-x}$,\quad $f_y=2e^{2y-x}$.

\textbf{$P(0,0)$'da:} $2y-x=0$; $f=e^0=1$, $f_x=-1$, $f_y=2$.
\[ L(x,y)=1+(-1)(x-0)+2(y-0)=1-x+2y . \]
\textbf{$Q(1,2)$'de:} $2y-x=3$; $f=e^3$, $f_x=-e^3$, $f_y=2e^3$.
\[ L(x,y)=e^3-e^3(x-1)+2e^3(y-2)=e^3\big(1-x+1+2y-4\big)=e^3(2y-x-2). \]
\end{cevap}

\begin{soru}[Soru 9 \foto{IMG-20261105-WA0038}, \foto{IMG-20261105-WA0039}, \foto{IMG-20261108-WA0085}, \foto{IMG-20261108-WA0086} (Problem 5a)]
$f(x,y)=x^2-3xy+y^2$ fonksiyonunun $P(1,2)$ noktasındaki $L(x,y)$ lineerleştirmesini ve
$B=\{(x,y):|x-1|\le0{,}1,\ |y-2|\le0{,}1\}$ bölgesindeki hatanın bir üst sınırını bulunuz.
\end{soru}
\begin{cevap}[Çözüm (adım adım)]
\textbf{Adım 1:} $f(1,2)=1-6+4=-1$.

\textbf{Adım 2:} $f_x=2x-3y\Rightarrow f_x(1,2)=2-6=-4$;\quad $f_y=-3x+2y\Rightarrow f_y(1,2)=-3+4=1$.

\textbf{Adım 3:} $L(x,y)=-1-4(x-1)+1\cdot(y-2)=-1-4x+4+y-2=-4x+y+1$.

\textbf{Adım 4:} $f_{xx}=2$, $f_{xy}=-3$, $f_{yy}=2$. $M$ mutlak değerlerden seçilir:
$M=\max\{|2|,|-3|,|2|\}=3$.

\textbf{Adım 5:}
\[ |E(x,y)|\le\frac12\cdot3\,(0{,}1+0{,}1)^2=1{,}5\cdot0{,}04=0{,}06 . \]
(Tahtada ``$f_x=2x-3y=f_x(1,2)=-4$'' biçiminde ``$\Rightarrow$'' yerine ``$=$'' yazılmıştı; anlamı
$f_x(1,2)=-4$'tür.)
\end{cevap}

\begin{soru}[Soru 10 \foto{IMG-20261105-WA0040}, \foto{IMG-20261108-WA0086} (Problem 5b)]
$f(x,y)=1+y+x\cos y$ fonksiyonunun $P(0,0)$ noktasındaki lineerleştirmesini ve
$B=\{(x,y):|x|\le0{,}2,\ |y|\le0{,}2\}$ bölgesindeki hatanın bir üst sınırını bulunuz.
\end{soru}
\begin{cevap}[Çözüm (adım adım)]
\textbf{Adım 1:} $f(0,0)=1+0+0=1$.

\textbf{Adım 2:} $f_x=\cos y\Rightarrow f_x(0,0)=1$;\quad $f_y=1-x\sin y\Rightarrow f_y(0,0)=1$.

\textbf{Adım 3:} $L(x,y)=1+1\cdot x+1\cdot y=1+x+y$.

\textbf{Adım 4 -- İkinci türevler:} $f_{xx}=0$, $f_{xy}=(\cos y)_y=-\sin y$, $f_{yy}=(1-x\sin y)_y=-x\cos y$.

\textbf{Adım 5 -- $B$ üzerinde üst sınır:} $|\sin y|\le|y|$ ve $|\cos y|\le1$ eşitsizlikleriyle
\[
|f_{xy}|=|\sin y|\le|y|\le0{,}2,\qquad |f_{yy}|=|x|\,|\cos y|\le0{,}2\cdot1=0{,}2,\qquad |f_{xx}|=0
\ \Rightarrow\ M=0{,}2 .
\]
\textbf{Adım 6:}
\[ |E(x,y)|\le\frac12\cdot0{,}2\,(0{,}2+0{,}2)^2=0{,}1\cdot0{,}16=0{,}016 . \]
\textbf{Dikkat:} İkinci türevler yalnızca $P(0,0)$'da hesaplansaydı hepsi $0$ olurdu; $M$, $B$'nin \emph{tamamındaki}
üst sınır olmalıdır.
\end{cevap}

%-------------------------------------------------------------------------------------------
\subsection{Tam diferansiyel ve yüksek mertebeden diferansiyeller (08.11 dersi)}

\begin{soru}[Soru 11 \foto{IMG-20261108-WA0065}, \foto{IMG-20261108-WA0066}, \foto{IMG-20261108-WA0067}]
Yarıçapı $r=3$ cm, yüksekliği $h=10$ cm olan bir dik dairesel silindirin yarıçapı $0{,}01$ cm arttırılıp yüksekliği
$0{,}1$ cm azaltılırsa hacmindeki değişim nedir?
\end{soru}
\begin{cevap}[Çözüm (adım adım)]
\textbf{Adım 1:} $V=\pi r^2h\Rightarrow dV=V_r\,dr+V_h\,dh=2\pi rh\,dr+\pi r^2\,dh$.

\textbf{Adım 2:} $r=3$, $h=10$, $dr=0{,}01$, $dh=-0{,}1$.

\textbf{Adım 3:}
\[
dV=2\pi\cdot3\cdot10\,(0{,}01)+\pi\cdot9\,(-0{,}1)=0{,}6\pi-0{,}9\pi=-0{,}3\pi\ \text{cm}^3\approx-0{,}94\ \text{cm}^3 .
\]
\textbf{Adım 4:} $V=\pi\cdot9\cdot10=90\pi$ $\Rightarrow$ yeni hacim $\approx90\pi-0{,}3\pi=89{,}7\pi$ cm$^3\approx281{,}8$ cm$^3$.

(Kontrol: gerçek değer $\pi\cdot3{,}01^2\cdot9{,}9\approx89{,}695\pi$ cm$^3$.)
\end{cevap}
\begin{dikkat}[Tahtada düzeltilen hata]
Tahtada son satır ``$90\pi-0{,}3\pi=89{,}7$ cm$^3$'' olarak yazılmıştı; $\pi$ düşmüştür. Doğrusu $89{,}7\pi$ cm$^3$
($\approx281{,}8$ cm$^3$).
\end{dikkat}

\begin{soru}[Soru 12 \foto{IMG-20261108-WA0068}, \foto{IMG-20261108-WA0069}]
$z=x^2y$ fonksiyonunun ikinci mertebeden diferansiyelini ($d^2z$) bulunuz.
\end{soru}
\begin{cevap}[Çözüm (adım adım)]
\textbf{Formül:} $dx$, $dy$ sabit artımlar kabul edilerek $d^2z=d(dz)=z_{xx}\,dx^2+2z_{xy}\,dx\,dy+z_{yy}\,dy^2$
(ikinci türevler sürekli olduğundan $z_{xy}=z_{yx}$).

\textbf{Türevler:} $z_x=2xy$, $z_y=x^2$; $z_{xx}=2y$, $z_{xy}=2x$, $z_{yy}=0$.
\[ d^2z=2y\,dx^2+2\cdot2x\,dx\,dy+0=2y\,dx^2+4x\,dx\,dy . \]
\end{cevap}

%-------------------------------------------------------------------------------------------
\subsection{Ek alıştırmalar}

\begin{soru}[\ekal Ek alıştırma 1 (tanımla kısmi türev)]
$f(x,y)=xy^2-x^3$ için $f_x(1,-1)$'i limit tanımıyla hesaplayınız ve kuralla doğrulayınız.
\end{soru}
\begin{cevap}
$f(1,-1)=1-1=0$; $f(1+h,-1)=(1+h)-(1+h)^3=1+h-1-3h-3h^2-h^3=-2h-3h^2-h^3$.
\[ f_x(1,-1)=\lim_{h\to0}\frac{-2h-3h^2-h^3}{h}=\lim_{h\to0}(-2-3h-h^2)=-2 . \]
Kuralla: $f_x=y^2-3x^2\Rightarrow f_x(1,-1)=1-3=-2$ $\checkmark$.
\end{cevap}

\begin{soru}[\ekal Ek alıştırma 2 (ikinci mertebe ve Clairaut)]
$f(x,y)=x\,e^{y}+y\ln x$ $(x>0)$ fonksiyonunun tüm ikinci mertebe kısmi türevlerini bulunuz; $f_{xy}=f_{yx}$
olduğunu gösteriniz.
\end{soru}
\begin{cevap}
$f_x=e^y+\dfrac yx$, \quad $f_y=xe^y+\ln x$.
\[
f_{xx}=-\frac{y}{x^2},\qquad f_{xy}=e^y+\frac1x,\qquad f_{yx}=e^y+\frac1x,\qquad f_{yy}=xe^y .
\]
$f_{xy}=f_{yx}$; karma türevler $x>0$ bölgesinde sürekli olduğundan bu Schwarz teoremiyle uyumludur.
\end{cevap}

\begin{soru}[\ekal Ek alıştırma 3 (teğet düzlem ve normal doğru)]
$z=3x^2-xy+y^2$ yüzeyinin $(1,2,5)$ noktasındaki teğet düzleminin ve normal doğrusunun denklemlerini bulunuz.
\end{soru}
\begin{cevap}
$f(1,2)=3-2+4=5$ (nokta yüzey üzerinde). $f_x=6x-y\Rightarrow f_x(1,2)=4$; $f_y=-x+2y\Rightarrow f_y(1,2)=3$.
\[ z-5=4(x-1)+3(y-2)\quad\Longrightarrow\quad z=4x+3y-5 . \]
Normal doğru (doğrultu $\langle4,3,-1\rangle$): $x=1+4t$, $y=2+3t$, $z=5-t$.
\end{cevap}

\begin{soru}[\ekal Ek alıştırma 4 (lineerleştirme ve hata sınırı)]
$f(x,y)=\ln x+\ln y$ fonksiyonunun $(1,1)$'deki lineerleştirmesini ve $B:\ |x-1|\le0{,}2,\ |y-1|\le0{,}2$ üzerindeki
hata sınırını bulunuz.
\end{soru}
\begin{cevap}
$f(1,1)=0$, $f_x=\dfrac1x\Rightarrow1$, $f_y=\dfrac1y\Rightarrow1$: \ $L(x,y)=(x-1)+(y-1)=x+y-2$.\\
$f_{xx}=-\dfrac1{x^2}$, $f_{yy}=-\dfrac1{y^2}$, $f_{xy}=0$. $B$'de $x,y\ge0{,}8$ olduğundan
$|f_{xx}|,|f_{yy}|\le\dfrac1{0{,}64}=1{,}5625=M$ (en büyük değer kutunun $0{,}8$ kenarında alınır).
\[ |E|\le\frac{1{,}5625}{2}\,(0{,}2+0{,}2)^2=0{,}78125\cdot0{,}16=0{,}125 . \]
\end{cevap}

\begin{soru}[\ekal Ek alıştırma 5 (birinci ve ikinci diferansiyel)]
$z=x\sin y$ için $dz$ ve $d^2z$'yi bulunuz.
\end{soru}
\begin{cevap}
$z_x=\sin y$, $z_y=x\cos y$ $\Rightarrow$ $dz=\sin y\,dx+x\cos y\,dy$.\\
$z_{xx}=0$, $z_{xy}=\cos y$, $z_{yy}=-x\sin y$ $\Rightarrow$ $d^2z=2\cos y\,dx\,dy-x\sin y\,dy^2$.
\end{cevap}

\begin{soru}[\ekal Ek alıştırma 6 (diferansiyelle yaklaşık değer)]
$(1{,}02)^{3{,}01}$ sayısını diferansiyel yardımıyla yaklaşık hesaplayınız.
\end{soru}
\begin{cevap}
$f(x,y)=x^y$, $(x_0,y_0)=(1,3)$, $dx=0{,}02$, $dy=0{,}01$. $f(1,3)=1$;
$f_x=y\,x^{y-1}\Rightarrow3$; $f_y=x^y\ln x\Rightarrow1\cdot\ln1=0$.
\[ df=3(0{,}02)+0(0{,}01)=0{,}06\quad\Rightarrow\quad(1{,}02)^{3{,}01}\approx1{,}06 . \]
(Hesap makinesi: $\approx1{,}0614$.)
\end{cevap}

\begin{soru}[\ekal Ek alıştırma 7 (yüzde değişim)]
Bir dirençte harcanan güç $P=\dfrac{V^2}{R}$'dir. Gerilim \%2, direnç \%1 artarsa güç yaklaşık yüzde kaç değişir?
\end{soru}
\begin{cevap}
$\ln P=2\ln V-\ln R\Rightarrow\dfrac{dP}{P}=2\dfrac{dV}{V}-\dfrac{dR}{R}=2(0{,}02)-0{,}01=0{,}03$.
Güç yaklaşık \textbf{\%3 artar}.
\end{cevap}

\begin{soru}[\ekal Ek alıştırma 8 (üç değişkenli ikinci diferansiyel)]
$w=xyz$ için $dw$ ve $d^2w$'yi bulunuz.
\end{soru}
\begin{cevap}
$w_x=yz$, $w_y=xz$, $w_z=xy$ $\Rightarrow$ $dw=yz\,dx+xz\,dy+xy\,dz$.\\
$w_{xx}=w_{yy}=w_{zz}=0$; $w_{xy}=z$, $w_{xz}=y$, $w_{yz}=x$ $\Rightarrow$
$d^2w=2z\,dx\,dy+2y\,dx\,dz+2x\,dy\,dz$.
\end{cevap}
